% !TEX TS-program = xelatex
\documentclass[a4paper]{article}
\usepackage[left=1in, right=1in, top=1in, bottom=1in]{geometry}
\usepackage{fontspec}
\usepackage{polyglossia}
\setdefaultlanguage{thai}
\setotherlanguage{english}
\usepackage{graphicx}

% Font setup for Thai with proper line breaking
\newfontfamily\thaifont{TH Sarabun New}[
  Script=Thai,
  Scale=1.0,
  WordSpace=1.0,
  LetterSpace=0.0
]
\newfontfamily\englishfont{TH Sarabun New}[
  Script=Latin,
  Scale=1.0
]
\setmainfont{TH Sarabun New}

% Enable Thai line breaking and hyphenation
\XeTeXlinebreaklocale "th"
\XeTeXlinebreakskip = 0pt plus 1pt minus 1pt

% Single spacing with proper line height for Thai
\usepackage{setspace}
\setstretch{1.2}

% 16pt body size with better line height
\makeatletter
\renewcommand\normalsize{%
  \@setfontsize\normalsize{16pt}{19.2pt}%
}
\normalsize
\makeatother

% Section titles formatting
\usepackage{titlesec}
\titleformat{\section}{\fontsize{18pt}{21.6pt}\selectfont\bfseries}{\thesection}{1em}{}
\titleformat{\subsection}{\fontsize{16pt}{19.2pt}\selectfont\bfseries}{\thesubsection}{1em}{}
\titleformat{\subsubsection}{\fontsize{16pt}{19.2pt}\selectfont\bfseries}{\thesubsubsection}{1em}{}

% TOC formatting
\usepackage{tocloft}
\renewcommand{\cftsecfont}{\fontsize{16pt}{19.2pt}\selectfont}
\renewcommand{\cftsubsecfont}{\fontsize{16pt}{19.2pt}\selectfont}
\renewcommand{\cftsecpagefont}{\fontsize{16pt}{19.2pt}\selectfont}
\renewcommand{\cftsubsecpagefont}{\fontsize{16pt}{19.2pt}\selectfont}

% Essential packages
\usepackage[hidelinks]{hyperref}
\usepackage{booktabs}
\usepackage{csquotes}

% Thai text formatting improvements for formal documents
\usepackage{ragged2e}
\justifying % Full justification for formal thesis
\hyphenpenalty=10000 % Disable hyphenation for Thai
\exhyphenpenalty=10000

% Bibliography setup for APA
\usepackage[style=apa,backend=biber]{biblatex}
\DeclareLanguageMapping{english}{english-apa}
\addbibresource{references.bib}

% Thai alphabet page numbering - complete implementation
\makeatletter
\def\@thaialph#1{%
  \ifcase#1\or ก\or ข\or ค\or ง\or จ\or ฉ\or ช\or ซ\or ฌ\or ญ\or
  ฎ\or ฏ\or ฐ\or ฑ\or ฒ\or ณ\or ด\or ต\or ถ\or ท\or ธ\or น\or บ\or ป\or ผ\or ฝ\or พ\or ฟ\or ภ\or ม\or
  ย\or ร\or ล\or ว\or ศ\or ษ\or ส\or ห\or ฬ\or อ\or ฮ\else\@ctrerr\fi
}
\def\thaialph#1{\expandafter\@thaialph\csname c@#1\endcsname}
\makeatother

% Formal Thai thesis paragraph settings
\setlength{\parindent}{2em}   % Standard Thai thesis indentation
\setlength{\parskip}{0pt}     % No space between paragraphs

% Better text justification for Thai
\tolerance=1000
\pretolerance=800
\emergencystretch=3em

% Cover page title setup
\date{}
\author{}
\title{
    \includegraphics[width=1in, height=1in, keepaspectratio]{logo.png}
    \\[2ex]
    {\fontsize{32pt}{36pt}\selectfont\textbf{Polyfunctional Robots}}
}

\begin{document}

% ----- Cover Page (no page number) -----
\maketitle
\thispagestyle{empty}

\vfill
\begin{center}
    {\fontsize{22pt}{26pt}\selectfont\textbf{
        นายคมชาญ วิเศษนคร
        \\ 663040419-1
    }}
\end{center}
\vfill

\vfill
\begin{center}
    {\fontsize{16pt}{20pt}\selectfont\textbf{
        รายงานฉบับนี้เป็นส่วนหนึ่งของรายวิชา EN813761 การสัมมนาทางวิศวกรรมคอมพิวเตอร์
        \\ สาขาวิชาวิศวกรรมคอมพิวเตอร์ มหาวิทยาลัยขอนแก่น
        \\ ภาคเรียนที่ 1 ปีการศึกษา 2568
    }}
\end{center}

\newpage

% ----- Start Thai letter page numbering for front matter -----
\pagenumbering{thaialph}
\setcounter{page}{1}

% ----- Abstract (Thai) -----
{\centering
    {\fontsize{18pt}{21.6pt}\selectfont\textbf{บทคัดย่อ}\par}
}
\vspace{1em}

% Thai abstract content goes here

\vspace{1em}

% ----- Abstract (English) -----
{\centering
    {\fontsize{18pt}{21.6pt}\selectfont\textbf{Abstract}\par}
}
\vspace{1em}

% English abstract content goes here

\vspace{1em}

\noindent
\textbf{คำสำคัญ:} หุ่นยนต์อเนกประสงค์, หุ่นยนต์โมดูลาร์, การควบคุมแบบลำดับชั้น, เซ็นเซอร์หลายโหมด, แอคชูเอเตอร์ปรับความแข็งได้

\vspace{0.5em}

\noindent
\textbf{Keywords:} Polyfunctional Robots, Modular Robotics, Hierarchical Control, Multimodal Sensors, Variable Stiffness Actuators

\newpage

% ----- Table of Contents -----
\tableofcontents
\newpage

% ----- Switch to Arabic numbering for main content -----
\pagenumbering{arabic}
\setcounter{page}{1}

% ----- Unnumbered front matter sections -----
\section*{คำนำ}
% คำนำ content goes here

\newpage

\section*{องค์ประกอบของรายงาน}
% รายละเอียดองค์ประกอบตามเกณฑ์การประเมิน

\newpage

% ----- Main Content Sections (Numbered) -----
\section{บทนำ}
หุ่นยนต์อเนกประสงค์ (Polyfunctional Robots) หรือหุ่นยนต์อเนกฟังก์ชัน เป็นระบบหุ่นยนต์ขั้นสูงที่ออกแบบให้สามารถปฏิบัติภารกิจที่หลากหลายและซับซ้อนภายในระบบเดียว โดยไม่จำเป็นต้องมีการปรับเปลี่ยนฮาร์ดแวร์หลักอย่างมีนัยสำคัญ \parencite{liang2025decoding} ซึ่งแตกต่างจากหุ่นยนต์แบบดั้งเดิมที่มักออกแบบมาเพื่อปฏิบัติงานเฉพาะทางเพียงอย่างเดียว หุ่นยนต์อเนกประสงค์สามารถปรับเปลี่ยนฟังก์ชันการทำงานผ่านการปรับโครงสร้างทางกายภาพ (Physical Reconfiguration) การเขียนโปรแกรมควบคุมใหม่ (Software Reconfiguration) หรือการผสมผสานโมดูลต่างๆ เข้าด้วยกัน \parencite{post2023modular}

ในยุคของการปฏิวัติอุตสาหกรรม 4.0 และการพัฒนาปัญญาประดิษฐ์ ความต้องการหุ่นยนต์ที่มีความยืดหยุ่นและสามารถปรับตัวได้กับสภาพแวดล้อมที่เปลี่ยนแปลงเพิ่มขึ้นอย่างรวดเร็ว \parencite{mohammadi2023mobile} หุ่นยนต์อเนกประสงค์จึงกลายเป็นทางเลือกที่มีความสำคัญสำหรับอุตสาหกรรมต่างๆ ตั้งแต่การผลิตและการประกอบชิ้นส่วน ไปจนถึงการแพทย์และการสำรวจอวกาศ เนื่องจากสามารถลดต้นทุนการลงทุนและเพิ่มประสิทธิภาพการใช้งานผ่านการใช้ระบบเดียวสำหรับหลายงาน

แนวคิดของหุ่นยนต์อเนกประสงค์มีความเชื่อมโยงอย่างใกล้ชิดกับหุ่นยนต์โมดูลาร์ (Modular Robots) และหุ่นยนต์ที่ปรับโครงสร้างได้ด้วยตนเอง (Self-Reconfiguring Robots) \parencite{seo2019modular} อย่างไรก็ตาม หุ่นยนต์อเนกประสงค์มีจุดเน้นที่แตกต่างออกไป คือ การเน้นที่ความสามารถในการปฏิบัติงานหลากหลายประเภทมากกว่าการเปลี่ยนรูปร่างหรือโครงสร้าง ทำให้เหมาะสำหรับการประยุกต์ใช้ในสภาพแวดล้อมที่ต้องการความเชี่ยวชาญในหลายด้านพร้อมกัน

\subsection{วัตถุประสงค์ของการศึกษา}

การศึกษานี้มีวัตถุประสงค์หลักเพื่อวิเคราะห์และสังเคราะห์องค์ความรู้เกี่ยวกับหุ่นยนต์อเนกประสงค์ในมุมมองทางวิศวกรรม โดยมีจุดมุ่งหมายเฉพาะ ดังนี้

\textbf{1. วิเคราะห์สถาปัตยกรรมและการออกแบบ} เพื่อศึกษาหลักการออกแบบสถาปัตยกรรมแบบโมดูลาร์ที่ปรับโครงสร้างได้ (Modular Reconfigurable Architecture) และระบบควบคุมแบบลำดับชั้น (Hierarchical Control Systems) ที่เป็นพื้นฐานสำคัญของหุ่นยนต์อเนกประสงค์ \parencite{tassi2024multimodal}

\textbf{2. ศึกษาเทคโนโลยีหลัก} โดยเฉพาะการบูรณาการเซ็นเซอร์แบบหลายโหมด (Multimodal Sensor Integration) \parencite{yang2024body} แอคชูเอเตอร์ปรับความแข็งได้ (Variable Stiffness Actuators) และการประยุกต์ใช้ปัญญาประดิษฐ์แบบโมเดลพื้นฐาน (Foundation Models) ในการควบคุม

\textbf{3. วิเคราะห์การประยุกต์ใช้} ในภาคอุตสาหกรรมสำคัญ รวมถึงการผลิตอัตโนมัติ การแพทย์ การสำรวจอวกาศ และการกู้ภัยพิบัติ พร้อมทั้งประเมินประสิทธิภาพและความคุ้มค่าทางเศรษฐศาสตร์

\textbf{4. ระบุความท้าทายและข้อจำกัด} ทั้งในด้านเทคนิคและการนำไปใช้งานจริง รวมถึงประเด็นด้านความปลอดภัยและมาตรฐานสากล เช่น ISO 10218-1:2025 \parencite{iso2025robotics}

\subsection{ขอบเขตการศึกษา}

การศึกษานี้มุ่งเน้นหุ่นยนต์อเนกประสงค์ในบริบททางวิศวกรรม โดยครอบคลุมระบบที่มีความสามารถในการปฏิบัติงานหลากหลายผ่านกลไกต่างๆ ดังนี้

\textbf{ขอบเขตด้านเทคนิค} การศึกษาครอบคลุมหุ่นยนต์ที่สามารถปรับเปลี่ยนฟังก์ชันผ่าน (1) การปรับโครงสร้างทางกายภาพแบบโมดูลาร์ (2) การเปลี่ยนแปลงอัลกอริทึมควบคุมและซอฟต์แวร์ และ (3) การผสมผสานโมดูลฮาร์ดแวร์ที่แตกต่างกัน

\textbf{ขอบเขตด้านการประยุกต์ใช้} เน้นการใช้งานในสภาพแวดล้อมอุตสาหกรรมและการค้า รวมถึงการแพทย์ การสำรวจ และการบริการ โดยไม่รวมถึงหุ่นยนต์เฉพาะทางที่ไม่สามารถปรับเปลี่ยนฟังก์ชันได้

\textbf{ขอบเขตด้านเวลา} การศึกษาเน้นงานวิจัยและพัฒนาตั้งแต่ปี ค.ศ. 2015 ถึงปัจจุบัน โดยเฉพาะอย่างยิ่งความก้าวหน้าในช่วง 5 ปีล่าสุดที่มีการประยุกต์ใช้ปัญญาประดิษฐ์และเทคโนโลยีการเรียนรู้ของเครื่อง

\subsection{คำจำกัดความสำคัญ}

เพื่อความชัดเจนในการศึกษา จึงกำหนดคำจำกัดความของแนวคิดสำคัญ ดังนี้

\textbf{หุ่นยนต์อเนกประสงค์ (Polyfunctional Robots)} หมายถึง ระบบหุ่นยนต์ที่สามารถปฏิบัติงานที่หลากหลายและแตกต่างกันได้ภายในระบบเดียว โดยมีความสามารถในการปรับเปลี่ยนฟังก์ชันการทำงานตามความต้องการของงานแต่ละประเภท

\textbf{หุ่นยนต์โมดูลาร์ (Modular Robots)} หมายถึง หุ่นยนต์ที่ประกอบด้วยโมดูลแยกส่วนที่สามารถเชื่อมต่อและแยกออกจากกันได้ เพื่อสร้างโครงสร้างและฟังก์ชันใหม่ตามต้องการ \parencite{bi2016survey}

\textbf{หุ่นยนต์ปรับโครงสร้างได้ (Self-Reconfiguring Robots)} หมายถึง หุ่นยนต์ที่สามารถเปลี่ยนแปลงรูปร่างและโครงสร้างของตนเองได้โดยอัตโนมัติ เพื่อให้เหมาะสมกับงานหรือสภาพแวดล้อมที่แตกต่างกัน \parencite{hameed2017modular}

\textbf{การควบคุมแบบลำดับชั้น (Hierarchical Control)} หมายถึง ระบบควบคุมที่จัดระดับการควบคุมเป็นชั้นๆ โดยชั้นบนมีหน้าที่วางแผนและตัดสินใจระดับสูง ส่วนชั้นล่างดำเนินการควบคุมรายละเอียดเฉพาะทาง

\subsection{ระเบียบวิธีการศึกษา}

การศึกษานี้ใช้วิธีการทบทวนวรรณกรรมเชิงพรรณนา (Descriptive Literature Review) ร่วมกับการวิเคราะห์เชิงเปรียบเทียบ โดยรวบรวมข้อมูลจากแหล่งข้อมูลทางวิชาการที่เชื่อถือได้ ประกอบด้วย

\textbf{แหล่งข้อมูลหลัก} วารสารวิชาการระดับนานาชาติที่ผ่านการประเมินโดยผู้ทรงคุณวุฒิ (Peer-reviewed Journals) เช่น IEEE Transactions on Robotics, International Journal of Robotics Research, และ Journal of Intelligent \& Robotic Systems

\textbf{แหล่งข้อมูลทุติยภูมิ} รายงานการประชุมวิชาการนานาชาติ (Conference Proceedings) มาตรฐานสากล และรายงานวิจัยจากสถาบันชั้นนำ เช่น NIST และ ISO

\textbf{เกณฑ์การคัดเลือกข้อมูล} เน้นงานวิจัยที่ตีพิมพ์ในช่วงปี ค.ศ. 2015-2025 มีการอ้างอิงและความน่าเชื่อถือสูง และเกี่ยวข้องโดยตรงกับหุ่นยนต์อเนกประสงค์หรือแนวคิดที่เกี่ยวข้อง

การวิเคราะห์ข้อมูลใช้การสังเคราะห์เชิงพรรณนา (Narrative Synthesis) โดยจัดกลุ่มข้อมูลตามประเด็นหลัก วิเคราะห์แนวโน้มและความสัมพันธ์ และสรุปเป็นองค์ความรู้ที่เป็นระบบ

\subsection{ประโยชน์ที่คาดว่าจะได้รับ}

การศึกษานี้คาดว่าจะให้ประโยชน์แก่ผู้เกี่ยวข้องหลายกลุ่ม ดังนี้

\textbf{สำหรับนักวิจัยและนักวิชาการ} เป็นการสังเคราะห์องค์ความรู้ที่เป็นปัจจุบันและครอบคลุม สามารถใช้เป็นฐานข้อมูลสำหรับการวิจัยต่อยอดในอนาคต

\textbf{สำหรับผู้ประกอบการและวิศวกร} ให้ข้อมูลสำคัญสำหรับการตัดสินใจลงทุนและการประยุกต์ใช้เทคโนโลยีหุ่นยนต์อเนกประสงค์ในภาคอุตสาหกรรม

\textbf{สำหรับนักศึกษาและผู้ที่สนใจ} เป็นแหล่งข้อมูลการเรียนรู้ที่เป็นระบบเกี่ยวกับเทคโนโลยีหุ่นยนต์ขั้นสูงและแนวโน้มการพัฒนาในอนาคต

\textbf{สำหรับหน่วยงานกำกับดูแล} ให้ข้อมูลประกอบการพิจารณาจัดทำนโยบายและมาตรฐานที่เกี่ยวข้องกับการใช้งานหุ่นยนต์อเนกประสงค์อย่างปลอดภัยและมีประสิทธิภาพ
\subsection{โครงสร้างของรายงาน}

รายงานฉบับนี้ได้จัดโครงสร้างเนื้อหาเป็น 6 บทหลัก และภาคผนวก เพื่อนำเสนอองค์ความรู้เกี่ยวกับหุ่นยนต์อเนกประสงค์อย่างเป็นระบบและครอบคลุม โดยมีรายละเอียดในแต่ละบทดังนี้

\textbf{บทที่ 1 บทนำ} นำเสนอความเป็นมาและความสำคัญของหุ่นยนต์อเนกประสงค์ในบริบทของการปฏิวัติอุตสาหกรรม 4.0 กำหนดวัตถุประสงค์และขอบเขตการศึกษา อธิบายคำจำกัดความของแนวคิดสำคัญ รวมถึงระเบียบวิธีการศึกษาและประโยชน์ที่คาดว่าจะได้รับจากการศึกษาครั้งนี้

\textbf{บทที่ 2 การทบทวนวรรณกรรมที่เกี่ยวข้อง} ทบทวนและวิเคราะห์เอกสารทางวิชาการที่เกี่ยวข้องกับหุ่นยนต์อเนกประสงค์ ครอบคลุมนิยามและแนวคิดพื้นฐาน สถาปัตยกรรมและการออกแบบระบบโมดูลาร์ เทคโนโลยีหลักและระบบควบคุมขั้นสูง รวมถึงการประยุกต์ใช้ในภาคอุตสาหกรรมต่างๆ เพื่อสร้างฐานความรู้สำหรับการวิเคราะห์และสังเคราะห์ในบทต่อไป

\textbf{บทที่ 3 การวิเคราะห์และอภิปราย} วิเคราะห์และเปรียบเทียบเทคโนโลยีหุ่นยนต์อเนกประสงค์ประเภทต่างๆ อภิปรายผลกระทบต่อภาคอุตสาหกรรมและสังคม ประเมินความท้าทายและข้อจำกัดของเทคโนโลยีปัจจุบัน รวมถึงการสังเคราะห์องค์ความรู้เพื่อเสนอกรอบแนวคิดและทิศทางการพัฒนาที่เหมาะสมกับบริบทของประเทศไทย

\textbf{บทที่ 4 เอกสารประมาณการงานฉบับสมบูรณ์} นำเสนอรายละเอียดการดำเนินงานที่สอดคล้องกับเกณฑ์การประเมินงานวิชาการ ระบุองค์ประกอบต่างๆ ของรายงาน แสดงความครบถ้วนของเนื้อหาทางวิชาการ การใช้แหล่งอ้างอิงที่เชื่อถือได้ และการประยุกต์ใช้เทคโนโลยี Generative AI อย่างเหมาะสมในกระบวนการวิจัย

\textbf{บทที่ 5 ความคิดสร้างสรรค์และความเรียบร้อย} แสดงความคิดสร้างสรรค์ในการนำเสนอข้อมูลผ่านการใช้แผนภูมิ ตาราง และรูปภาพประกอบที่ช่วยให้เข้าใจเนื้อหาได้ง่ายขึ้น รวมถึงการจัดรูปแบบเอกสารที่เป็นระเบียบและสวยงาม เพื่อเพิ่มประสิทธิภาพในการสื่อสารองค์ความรู้ทางวิชาการ

\textbf{บทที่ 6 สรุป} สรุปผลการศึกษาและการค้นพบที่สำคัญ เสนอข้อเสนอแนะสำหรับการนำไปประยุกต์ใช้ในภาคปฏิบัติ ระบุทิศทางการวิจัยและพัฒนาในอนาคต รวมถึงการให้ข้อเสนอแนะเชิงนโยบายสำหรับการส่งเสริมและพัฒนาเทคโนโลยีหุ่นยนต์อเนกประสงค์ในประเทศไทย

\textbf{ภาคผนวก} ประกอบด้วยคำศัพท์เทคนิคที่สำคัญพร้อมคำอธิบายเป็นภาษาไทยและภาษาอังกฤษ ตารางเปรียบเทียบเทคโนโลยีต่างๆ และข้อมูลเสริมอื่นๆ ที่สนับสนุนเนื้อหาหลักของรายงาน เพื่อให้ผู้อ่านสามารถศึกษารายละเอียดเพิ่มเติมได้ตามความสนใจ

การจัดโครงสร้างเนื้อหาดังกล่าวมีจุดมุ่งหมายเพื่อให้ผู้อ่านสามารถติดตามเนื้อหาได้อย่างต่อเนื่องและเป็นระบบ โดยเริ่มจากการสร้างความเข้าใจพื้นฐาน ผ่านการทบทวนองค์ความรู้ที่มีอยู่ จนถึงการวิเคราะห์เชิงลึกและการสรุปผลที่สามารถนำไปประยุกต์ใช้ได้จริง

\section{การทบทวนวรรณกรรมที่เกี่ยวข้อง}
การทบทวนวรรณกรรมในบทนี้มุ่งเน้นการสังเคราะห์องค์ความรู้ที่เกี่ยวข้องกับหุ่นยนต์อเนกประสงค์จากแหล่งข้อมูลทางวิชาการที่เชื่อถือได้ โดยครอบคลุมตั้งแต่แนวคิดพื้นฐานไปจนถึงเทคโนโลยีขั้นสูงและการประยุกต์ใช้ในภาคปฏิบัติ เพื่อสร้างฐานความรู้ที่แข็งแกร่งสำหรับการวิเคราะห์และการพัฒนาเทคโนโลยีต่อไป

\subsection{แนวคิดและนิยามหุ่นยนต์อเนกประสงค์}

แนวคิดของหุ่นยนต์อเนกประสงค์ได้รับการพัฒนาขึ้นจากความต้องการในการสร้างระบบหุ่นยนต์ที่สามารถปรับตัวเข้ากับงานที่หลากหลายได้อย่างมีประสิทธิภาพ โดยไม่จำเป็นต้องออกแบบหุ่นยนต์ใหม่สำหรับแต่ละงาน \parencite{liang2025decoding} ในวรรณกรรมทางวิชาการพบว่ามีการใช้คำศัพท์ที่หลากหลายเพื่ออธิบายแนวคิดที่คล้ายกัน ได้แก่ "หุ่นยนต์หลายฟังก์ชัน (Multifunctional Robots)" "หุ่นยนต์อเนกประสงค์ (Versatile Robots)" และ "หุ่นยนต์ทั่วไป (Generalist Robots)" ซึ่งแต่ละคำศัพท์มีนัยและขอบเขตที่แตกต่างกันเล็กน้อย

การศึกษาเชิงสำรวจของ \textcite{bi2016survey} ได้แสดงให้เห็นว่าหุ่นยนต์อเนกประสงค์มีความเชื่อมโยงอย่างใกล้ชิดกับหุ่นยนต์โมดูลาร์ที่สามารถปรับโครงสร้างได้ โดยทั้งสองแนวทางมีเป้าหมายร่วมกันในการสร้างระบบที่มีความยืดหยุ่นและสามารถปรับตัวได้ อย่างไรก็ตาม หุ่นยนต์อเนกประสงค์เน้นที่ความสามารถในการปฏิบัติงานหลากหลายมากกว่าการเปลี่ยนแปลงโครงสร้างทางกายภาพ

งานวิจัยของ \textcite{hameed2017modular} ได้จำแนกหุ่นยนต์อเนกประสงค์ออกเป็น 3 ประเภทหลัก ได้แก่ หุ่นยนต์แบบเปลี่ยนโครงสร้างได้ (Reconfigurable) หุ่นยนต์แบบโมดูลาร์ (Modular) และหุ่นยนต์แบบปรับตัวได้ (Adaptive) โดยแต่ละประเภทมีกลไกและวิธีการในการบรรลุความสามารถอเนกประสงค์ที่แตกต่างกัน การจำแนกนี้ช่วยให้เข้าใจขอบเขตและข้อจำกัดของเทคโนโลยีแต่ละประเภทได้ชัดเจนยิ่งขึ้น

\textcite{seo2019modular} ได้เสนอกรอบแนวคิดที่ครอบคลุมสำหรับการออกแบบหุ่นยนต์อเนกประสงค์ โดยเน้นที่การบูรณาการระหว่างฮาร์ดแวร์และซอฟต์แวร์ การศึกษานี้แสดงให้เห็นว่าการบรรลุความสามารถอเนกประสงค์ต้องอาศัยการประสานงานระหว่างการออกแบบเชิงกล ระบบควบคุม และอัลกอริทึมการเรียนรู้อย่างเป็นระบบ

ความก้าวหน้าล่าสุดในการใช้โมเดลพื้นฐาน (Foundation Models) ได้เปิดมิติใหม่ให้กับแนวคิดหุ่นยนต์อเนกประสงค์ \textcite{wang2023robot} ได้ทบทวนการประยุกต์ใช้โมเดลภาษาขนาดใหญ่และโมเดลพื้นฐานในการสร้างนโยบายการทำงานแบบทั่วไป (Generalist Policies) ที่สามารถปรับใช้กับงานใหม่ได้โดยต้องการข้อมูลฝึกเพิ่มเติมเพียงเล็กน้อย ซึ่งเป็นการเปลี่ยนแปลงสำคัญจากแนวทางแบบดั้งเดิมที่ต้องเขียนโปรแกรมเฉพาะสำหรับแต่ละงาน

\subsection{สถาปัตยกรรมและการออกแบบระบบ}

สถาปัตยกรรมของหุ่นยนต์อเนกประสงค์เป็นองค์ประกอบสำคัญที่กำหนดความสามารถและข้อจำกัดของระบบ การศึกษาของ \textcite{post2023modular} ได้นำเสนอการสำรวจเชิงลึกเกี่ยวกับสถาปัตยกรรมของหุ่นยนต์โมดูลาร์ที่สามารถปรับโครงสร้างได้ด้วยตนเอง โดยครอบคลุมกลไกการเชื่อมต่อ (Docking Mechanisms) โมดูลฟังก์ชัน (Functional Modules) และความสามารถในการปรับโครงสร้างจากมุมมองทางวิศวกรรมเครื่องกล

การออกแบบกลไกการเชื่อมต่อมีความสำคัญอย่างยิ่งต่อความสำเร็จของหุ่นยนต์อเนกประสงค์ งานวิจัยของ \textcite{liang2025decoding} ได้วิเคราะห์หลักการออกแบบกลไกการเชื่อมต่อสำหรับหุ่นยนต์โมดูลาร์ที่สามารถปรับโครงสร้างได้ โดยครอบคลุมทั้งการออกแบบแบบกำหนดลำดับ (Deterministic) และแบบสุ่ม (Stochastic) การศึกษานี้แสดงให้เห็นถึงความซับซ้อนในการออกแบบระบบที่ต้องรองรับการเปลี่ยนแปลงโครงสร้างอย่างต่อเนื่องในขณะเดียวกันต้องรักษาความแข็งแรงเชิงกลและความเสถียรของระบบ

\textcite{krishnan2023deep} ได้เสนอแนวทางการใช้การเรียนรู้เชิงลึกในการออกแบบกลไกหุ่นยนต์ โดยเฉพาะสำหรับสถาปัตยกรรมแบบหลายฟังก์ชัน การศึกษานี้แสดงให้เห็นว่าการใช้เครือข่ายประสาทเทียมสามารถช่วยให้การออกแบบระบบที่ซับซ้อนเป็นไปได้มากขึ้น โดยเฉพาะในการหาจุดสมดุลระหว่างข้อกำหนดทางวิศวกรรมที่หลากหลายและขัดแย้งกัน

การพัฒนาสถาปัตยกรรมแบบต่อเนื่อง (Continuum Architecture) เป็นอีกทิศทางหนึ่งที่น่าสนใจสำหรับหุ่นยนต์อเนกประสงค์ \textcite{ieee2023dynamic} ได้นำเสนอการจำลองพลวัตของหุ่นยนต์ต่อเนื่องที่ขับเคลื่อนด้วยแอคชูเอเตอร์แบบ bellows โดยใช้สูตรออยเลอร์-ลากรานจ์ หุ่นยนต์ประเภทนี้มีความสามารถในการปรับเปลี่ยนรูปร่างและคุณสมบัติทางกลได้อย่างต่อเนื่อง ทำให้เหมาะสำหรับงานที่ต้องการความยืดหยุ่นสูง

การประยุกต์ใช้การคำนวณเชิงสัณฐาน (Morphological Computation) ในการออกแบบหุ่นยนต์อเนกประสงค์เป็นแนวทางที่ได้รับความสนใจเพิ่มขึ้น \textcite{sitti2021morphological} ได้ศึกษาการประยุกต์ใช้กลไกการกระจายเมล็ดพันธุ์พืชในการออกแบบหุ่นยนต์อ่อนรุ่นใหม่ที่สามารถเคลื่อนที่ได้หลายรูปแบบ เช่น การบิน การเคลื่อนคลาน และการเจาะผ่านคุณสมบัติโครงสร้างแบบพาสซีฟ

\textcite{li2025finite} ได้นำเสนอการวิเคราะห์องค์ประกอบไฟไนต์และการออกแบบเพิ่มประสิทธิภาพเชิงโครงสร้างสำหรับแขนหุ่นยนต์หลายฟังก์ชันสำหรับรถขยะ การศึกษานี้แสดงให้เห็นถึงการประยุกต์ใช้เครื่องมือการวิเคราะห์ขั้นสูงในการออกแบบระบบที่ต้องรองรับการทำงานหลากหลาย โดยสามารถลดน้ำหนักลงได้ 14.28% ในขณะที่รักษาความแข็งแรงและความแกร่งตามข้อกำหนด

\subsection{เทคโนโลยีหลักและระบบควบคุม}

ระบบควบคุมเป็นหัวใจสำคัญที่ทำให้หุ่นยนต์อเนกประสงค์สามารถปรับเปลี่ยนระหว่างงานต่างๆ ได้อย่างราบรื่นและมีประสิทธิภาพ \textcite{tassi2024multimodal} ได้นำเสนอกรอบการทำงานของการควบคุมแบบหลายโหมดและปรับตัวได้ผ่านการเขียนโปรแกรมกำลังสองแบบลำดับชั้น (Hierarchical Quadratic Programming) ซึ่งช่วยให้หุ่นยนต์สามารถปรับเปลี่ยนพฤติกรรมการตอบสนองได้ตามลักษณะของงานและสภาพแวดล้อม

การพัฒนาเซ็นเซอร์แบบหลายโหมดเป็นเทคโนโลยีสำคัญที่ขับเคลื่อนความสามารถอเนกประสงค์ของหุ่นยนต์ \textcite{yang2024body} ได้นำเสนอการออกแบบผิวหนังหุ่นยนต์ขนาดตัวที่ใช้โมดูลการรับรู้แบบหลายโหมดแบบกระจาย ระบบนี้สามารถรับรู้ข้อมูลการสัมผัส แรงกด อุณหภูมิ และการเคลื่อนไหวได้พร้อมกันผ่านเซ็นเซอร์ที่กระจายอยู่ทั่วผิวหน้าของหุ่นยนต์ ทำให้สามารถปรับพฤติกรรมได้ตามลักษณะของวัตถุและงานที่แตกต่างกัน

\textcite{athar2023vistac} ได้พัฒนาเซ็นเซอร์แบบรวม (Unified Multimodal Sensing) ที่ผสมผสานการมองเห็นและการสัมผัสสำหรับการจับจัดหุ่นยนต์ เทคโนโลยี VisTac นี้ช่วยให้หุ่นยนต์สามารถตัดสินใจเลือกกลยุทธ์การจับยึดที่เหมาะสมกับวัตถุแต่ละประเภทได้อย่างอัตโนมัติ ซึ่งเป็นความสามารถสำคัญสำหรับหุ่นยนต์อเนกประสงค์ที่ต้องจัดการกับวัตถุหลากหลายประเภท

แอคชูเอเตอร์ปรับความแข็งได้เป็นอีกหนึ่งเทคโนโลยีหลักที่ช่วยให้หุ่นยนต์อเนกประสงค์สามารถปรับพฤติกรรมทางกลได้ตามลักษณะของงาน \textcite{ieee2024compact} ได้นำเสนอแอคชูเอเตอร์ปรับความแข็งแบบกะทัดรัดสำหรับการเคลื่อนที่แบบคล่องตัว ระบบนี้สามารถปรับระดับความแข็งแรงได้ตั้งแต่การเคลื่อนไหวที่นุ่มนวลสำหรับการปฏิสัมพันธ์กับมนุษย์ไปจนถึงการเคลื่อนไหวที่แข็งแรงสำหรับงานที่ต้องการความแม่นยำสูง

\textcite{jin2021origami} ได้พัฒนาแอคชูเอเตอร์อ่อนที่ได้แรงบันดาลใจจากศิลปะโอริกามิสำหรับการรับรู้สิ่งเร้าและการประยุกต์ใช้ในหุ่นยนต์เคลื่อนคลาน แอคชูเอเตอร์ประเภทนี้สามารถทำหน้าที่ทั้งเป็นเซ็นเซอร์และตัวขับเคลื่อนในเวลาเดียวกัน ซึ่งช่วยลดความซับซ้อนของระบบและเพิ่มความสามารถในการปรับตัว

การบูรณาการระบบควบคุมแบบลำดับชั้นเป็นกลยุทธ์สำคัญในการจัดการความซับซ้อนของหุ่นยนต์อเนกประสงค์ \textcite{ieee2019continuous} ได้นำเสนอแนวทางการเปลี่ยนงานแบบต่อเนื่องสำหรับตัวควบคุมหุ่นยนต์ที่ใช้การเขียนโปรแกรมกำลังสองแบบลำดับชั้น ระบบนี้ช่วยให้หุ่นยนต์สามารถเพิ่ม ลบ และเปลี่ยนงานได้โดยไม่เกิดการหยุดชะงักหรือการเปลี่ยนแปลงอย่างกะทันหัน

การประยุกต์ใช้การเรียนรู้ของเครื่องแบบต่อเนื่องเป็นแนวทางที่มีศักยภาพสูงสำหรับหุ่นยนต์อเนกประสงค์ \textcite{thuruthel2021survey} ได้ทบทวนการประยุกต์ใช้การเรียนรู้ของเครื่องสำหรับการควบคุมหุ่นยนต์ต่อเนื่อง โดยเน้นที่ความสามารถในการเรียนรู้และปรับตัวกับงานใหม่โดยไม่สูญเสียความรู้เดิม ซึ่งเป็นคุณสมบัติสำคัญสำหรับระบบที่ต้องปฏิบัติงานหลากหลายประเภท

\subsection{การประยุกต์ใช้ในภาคอุตสาหกรรม}

การประยุกต์ใช้หุ่นยนต์อเนกประสงค์ในภาคอุตสาหกรรมได้รับความสนใจเพิ่มขึ้นอย่างรวดเร็ว โดยเฉพาะในบริบทของการปฏิวัติอุตสาหกรรม 4.0 \textcite{mohammadi2023mobile} ได้นำเสนอการทบทวนเชิงลึกเกี่ยวกับโมบายแมนิปิวเลเตอร์ในการประยุกต์ใช้อุตสาหกรรม 4.0 การศึกษานี้แสดงให้เห็นว่าหุ่นยนต์เคลื่อนที่ที่มีแขนจับได้กลายเป็นองค์ประกอบสำคัญของระบบการผลิตที่ชาญฉลาดและยืดหยุ่น

ในภาคการผลิต การบูรณาการระหว่างหุ่นยนต์เคลื่อนที่และหุ่นยนต์ร่วมมือ (Collaborative Robots) ได้รับการพัฒนาเป็นแนวทางใหม่ที่มีศักยภาพสูง \textcite{vitolo2022mobile} ได้นำเสนอการออกแบบเบื้องต้นของอินเทอร์เฟซทางเมคาทรอนิกส์สำหรับการบูรณาการหุ่นยนต์เคลื่อนที่และโคบอทโดยใช้แนวทาง Model-Based Systems Engineering (MBSE) การบูรณาการนี้ช่วยให้สามารถสร้างระบบการผลิตที่มีความยืดหยุ่นสูงและสามารถปรับเปลี่ยนประเภทผลิตภัณฑ์ได้อย่างรวดเร็ว

การวัดและประเมินประสิทธิภาพของหุ่นยนต์อเนกประสงค์ในสภาพแวดล้อมอุตสาหกรรมเป็นประเด็นสำคัญที่ต้องการมาตรฐานที่ชัดเจน \textcite{bostelman2016mobile} จากสถาบัน NIST (National Institute of Standards and Technology) ได้พัฒนามาตรฐานการวัดประสิทธิภาพของโมบายแมนิปิวเลเตอร์สำหรับงานประกอบเพื่อการผลิตการศึกษานี้ได้กำหนดเมตริกและวิธีการทดสอบที่มาตรฐานสำหรับการประเมินความสามารถของหุ่นยนต์อเนกประสงค์ในการปฏิบัติงานจริง

\textcite{xie2024pose} ได้นำเสนอการเพิ่มประสิทธิภาพท่าทางสำหรับการจับยึดของโมบายแมนิปิวเลเตอร์โดยใช้เกณฑ์การจัดการแบบไฮบริด การศึกษานี้แสดงให้เห็นถึงความสำคัญของการวางตำแหน่งฐานเคลื่อนที่ที่เหมาะสมเพื่อให้ได้ประสิทธิภาพการจับยึดที่ดีที่สุด ซึ่งเป็นปัจจัยสำคัญสำหรับหุ่นยนต์อเนกประสงค์ที่ต้องจัดการกับวัตถุหลากหลายประเภท

ในภาคการดูแลสุขภาพ การพัฒนาหุ่นยนต์อเนกประสงค์ได้รับความสนใจเป็นพิเศษเนื่องจากความต้องการที่หลากหลายในสภาพแวดล้อมการดูแลผู้ป่วย \textcite{stueckler2023hollie} ได้นำเสนอหุ่นยนต์เคลื่อนที่แบบสองแขนหลายฟังก์ชันที่รองรับการประยุกต์ใช้การดูแลที่หลากหลาย ระบบนี้สามารถปรับเปลี่ยนระหว่างการช่วยเหลือในกิจวัตรประจำวัน การขนส่งยาและอุปกรณ์การแพทย์ และการให้การสนับสนุนทางอารมณ์ได้ตามความต้องการของผู้ป่วยแต่ละคน

การพัฒนาหุ่นยนต์อเนกประสงค์สำหรับการประยุกต์ใช้ในสภาพแวดล้อมที่รุนแรงหรือเป็นอันตราย เช่น การสำรวจอวกาศหรือการกู้ภัยพิบัติ ต้องการความทนทานและความสามารถในการปรับตัวเป็นพิเศษ หุ่นยนต์เหล่านี้ต้องสามารถปรับเปลี่ยนหน้าที่ได้อย่างรวดเร็วเมื่อสถานการณ์เปลี่ยนแปลง เช่น จากการสำรวจเป็นการกู้ภัย หรือจากการรวบรวมข้อมูลเป็นการซ่อมแซมอุปกรณ์

การพัฒนาระบบควบคุมหลายหุ่นยนต์ที่ประสานงานกันได้กลายเป็นแนวโน้มสำคัญในการประยุกต์ใช้อุตสาหกรรม \textcite{ieee2015prioritized} ได้นำเสนออัลกอริทึมการวางแผนแบบจัดลำดับความสำคัญสำหรับการประสานวิถีการเคลื่อนที่ของหุ่นยนต์เคลื่อนที่หลายตัว ซึ่งช่วยให้ระบบหุ่นยนต์อเนกประสงค์หลายตัวสามารถทำงานร่วมกันได้อย่างมีประสิทธิภาพในสภาพแวดล้อมที่ซับซ้อน

\subsection{ช่องว่างความรู้และแนวทางการวิจัย}

จากการทบทวนวรรณกรรมที่กว้างขวางพบว่ามีช่องว่างความรู้สำคัญหลายประการที่ต้องการการวิจัยและพัฒนาต่อไป แม้ว่าจะมีความก้าวหน้าอย่างมากในแต่ละองค์ประกอบของหุ่นยนต์อเนกประสงค์ แต่การบูรณาการและการประยุกต์ใช้ในสถานการณ์จริงยังคงมีความท้าทายที่สำคัญ

\textbf{ข้อจำกัดด้านการทดสอบในสภาพแวดล้อมจริง} การศึกษาส่วนใหญ่ยังคงเป็นการทดสอบในสภาพแวดล้อมที่ควบคุมได้หรือการจำลองด้วยคอมพิวเตอร์ การทดสอบระบบหุ่นยนต์อเนกประสงค์ที่มีความสามารถสูงในสภาพแวดล้อมจริงที่มีความไม่แน่นอนและความซับซ้อนสูงยังมีอยู่จำกัด ซึ่งเป็นอุปสรรคสำคัญต่อการนำไปประยุกต์ใช้ในภาคอุตสาหกรรมและการค้าจริง

\textbf{การขาดแคลนกรอบแนวคิดแบบบูรณาการ} แม้ว่าจะมีการพัฒนาเทคโนโลยีต่างๆ ที่เกี่ยวข้องกับหุ่นยนต์อเนกประสงค์อย่างต่อเนื่อง แต่ยังขาดกรอบแนวคิดที่เชื่อมโยงการปรับโครงสร้างทางกลกับพฤติกรรมที่เรียนรู้ได้อย่างเป็นระบบ การพัฒนาส่วนใหญ่ยังคงแยกส่วนระหว่างการออกแบบฮาร์ดแวร์และการพัฒนาซอฟต์แวร์ ซึ่งไม่เอื้อต่อการสร้างระบบที่มีประสิทธิภาพสูงสุด

\textbf{ประเด็นด้านประสิทธิภาพการใช้พลังงาน} การศึกษาเกี่ยวกับการจัดการพลังงานในระบบหุ่นยนต์อเนกประสงค์ยังมีอยู่อย่างจำกัด โดยเฉพาะในระบบที่ต้องการการปรับเปลี่ยนบ่อยครั้งหรือการทำงานต่อเนื่องเป็นเวลานาน การปรับโครงสร้างและการเปลี่ยนโหมดการทำงานมักต้องใช้พลังงานเพิ่มเติม ซึ่งเป็นปัจจัยสำคัญที่ส่งผลต่อความเป็นไปได้ในการประยุกต์ใช้จริง

\textbf{มาตรฐานการประเมินประสิทธิภาพที่ขาดความเป็นเอกภาพ} ปัจจุบันยังไม่มีมาตรฐานสากลที่ยอมรับร่วมกันสำหรับการประเมินประสิทธิภาพของหุ่นยนต์อเนกประสงค์ในมิติต่างๆ เช่น ความยืดหยุ่น ความน่าเชื่อถือ เวลาในการปรับเปลี่ยนงาน และความแม่นยำในการปฏิบัติงานที่หลากหลาย ความขาดแคลนนี้ทำให้การเปรียบเทียบและการประเมินผลงานวิจัยต่างๆ เป็นไปได้อย่างจำกัด

\textbf{ความท้าทายด้านความปลอดภัยและความเชื่อถือได้} การทำงานในสภาพแวดล้อมที่มีมนุษย์ร่วมด้วยต้องการระดับความปลอดภัยและความเชื่อถือได้สูง โดยเฉพาะเมื่อหุ่นยนต์สามารถเปลี่ยนพฤติกรรมหรือโครงสร้างได้ การพัฒนาระบบการรับประกันความปลอดภัยที่ครอบคลุมการเปลี่ยนแปลงแบบไดนามิกและการทำงานหลายโหมดยังคงเป็นความท้าทายที่สำคัญ

\textbf{การขาดแคลนข้อมูลสำหรับการเรียนรู้ข้ามโดเมน} แม้ว่าการใช้โมเดลพื้นฐานจะแสดงศักยภาพที่น่าสนใจ แต่ข้อมูลฝึกสำหรับหุ่นยนต์อเนกประสงค์ที่ครอบคลุมงานหลากหลายประเภทยังมีอยู่จำกัด โดยเฉพาะข้อมูลที่มีคุณภาพสูงสำหรับการเรียนรู้การปรับเปลี่ยนระหว่างงานที่แตกต่างกันอย่างมีนัยสำคัญ

\textbf{แนวทางการวิจัยในอนาคต} จากการวิเคราะห์ช่องว่างความรู้ดังกล่าว สามารถระบุแนวทางการวิจัยที่มีความสำคัญสูงได้ดังนี้

การพัฒนาแพลตฟอร์มทดสอบมาตรฐานสำหรับหุ่นยนต์อเนกประสงค์ในสภาพแวดล้อมที่หลากหลาย เพื่อให้สามารถประเมินประสิทธิภาพได้อย่างเป็นระบบและเปรียบเทียบได้ การพัฒนาแพลตฟอร์มนี้ควรครอบคลุมทั้งการทดสอบในสภาพแวดล้อมที่ควบคุมได้และการทดสอบในสถานการณ์จริงที่มีความซับซ้อน

การสร้างกรอบแนวคิดแบบบูรณาการที่เชื่อมโยงการออกแบบฮาร์ดแวร์ การพัฒนาซอฟต์แวร์ และการเรียนรู้ของเครื่องเข้าด้วยกัน โดยเฉพาะการพัฒนาวิธีการที่ช่วยให้การปรับโครงสร้างทางกายภาพและการเรียนรู้พฤติกรรมสามารถเกิดขึ้นพร้อมกันและส่งเสริมซึ่งกันและกัน

การวิจัยเชิงลึกเกี่ยวกับการจัดการพลังงานและการเพิ่มประสิทธิภาพการใช้พลังงานในระบบหุ่นยนต์อเนกประสงค์ รวมถึงการพัฒนาอัลกอริทึมการตัดสินใจที่คำนึงถึงการใช้พลังงานในการเลือกวิธีการปฏิบัติงาน

การพัฒนาชุดข้อมูลขนาดใหญ่และหลากหลายสำหรับการฝึกฝนหุ่นยนต์อเนกประสงค์ โดยเฉพาะข้อมูลที่เกี่ยวข้องกับการปรับเปลี่ยนระหว่างงานที่แตกต่างกันและการทำงานในสภาพแวดล้อมที่ไม่แน่นอน รวมถึงการพัฒนาเทคนิคการเรียนรู้ที่สามารถใช้ประโยชน์จากข้อมูลที่มีอยู่อย่างจำกัดได้อย่างมีประสิทธิภาพ

การสร้างมาตรฐานความปลอดภัยและแนวทางปฏิบัติที่ดีสำหรับหุ่นยนต์อเนกประสงค์ที่ทำงานในสภาพแวดล้อมที่มีมนุษย์ร่วมด้วย โดยเฉพาะการพัฒนาระบบการตรวจจับและการตอบสนองต่อสถานการณ์ที่ไม่คาดคิดระหว่างการเปลี่ยนโหมดการทำงาน

การพัฒนาเทคนิคการประเมินและการรับรองประสิทธิภาพที่ครอบคลุม รวมถึงการสร้างเมตริกใหม่ที่สามารถวัดความสามารถอเนกประสงค์ได้อย่างเป็นปริมาณ เช่น ดัชนีความยืดหยุ่น ค่าประสิทธิภาพการปรับตัว และตัวชี้วัดความคุ้มค่าทางเศรษฐศาสตร์

การศึกษาผลกระทบทางสังคมและเศรษฐกิจของการนำหุ่นยนต์อเนกประสงค์ไปใช้ในภาคต่างๆ รวมถึงการวิเคราะห์ความเป็นไปได้ในการปรับตัวของแรงงานและการเปลี่ยนแปลงของโครงสร้างอุตสาหกรรม

จากการทบทวนวรรณกรรมครั้งนี้แสดงให้เห็นว่าหุ่นยนต์อเนกประสงค์เป็นสาขาที่มีศักยภาพสูงและกำลังพัฒนาอย่างรวดเร็ว อย่างไรก็ตาม ยังคงมีความท้าทายและช่องว่างความรู้ที่สำคัญที่ต้องการการวิจัยอย่างเป็นระบบและการร่วมมือระหว่างสาขาวิชาต่างๆ เพื่อให้สามารถบรรลุเป้าหมายในการสร้างหุ่นยนต์ที่มีความสามารถอเนกประสงค์อย่างแท้จริงและนำไปประยุกต์ใช้ได้ในวงกว้าง

\section{ระเบียบวิธีวิจัย}
% EXTRACT from Section 1 + EXPAND
% 3.1 แนวทางการศึกษา
% 3.2 การรวบรวมข้อมูล  
% 3.3 เกณฑ์การคัดเลือกแหล่งข้อมูล
% 3.4 การวิเคราะห์และสังเคราะห์ข้อมูล
% 3.5 ข้อจำกัดของการศึกษา

\section{ผลการศึกษาและการวิเคราะห์}
% NEW SYNTHESIS CHAPTER - Critical analysis
% 4.1 การวิเคราะห์เชิงเปรียบเทียบเทคโนโลยี
% 4.2 การประเมินประสิทธิภาพและข้อจำกัด
% 4.3 แนวโน้มการพัฒนาและทิศทางอนาคต  
% 4.4 ความท้าทายและปัจจัยสำคัญ
% 4.5 กรอบแนวคิดเชิงบูรณาการ [SYNTHESIS]

\section{การอภิปรายผล}
% CRITICAL DISCUSSION - Higher-order analysis
% 5.1 การอภิปรายผลการวิเคราะห์
% 5.2 ผลกระทบต่อภาคอุตสาหกรรม
% 5.3 ข้อเสนอแนะเชิงนโยบายและการปฏิบัติ
% 5.4 ข้อจำกัดและข้อควรระวัง

\section{สรุปและข้อเสนอแนะ}
% COMPREHENSIVE CONCLUSIONS
% 6.1 สรุปผลการศึกษา
% 6.2 ข้อเสนอแนะสำหรับการวิจัยต่อไป  
% 6.3 ข้อเสนอแนะสำหรับการนำไปประยุกต์ใช้
% 6.4 คำสุดท้าย

% ----- Bibliography -----
\printbibliography[title=เอกสารอ้างอิง]

% ----- Appendices -----
\appendix

\section{ภาคผนวก ก: คำศัพท์เทคนิค}
% Technical glossary

\section{ภาคผนวก ข: ตารางเปรียบเทียบเทคโนโลยี}
% Technology comparison tables

\end{document}